\section{Complementary lower-bound and decision checks}
\label{sec:witness-comparison}

The block study extends an earlier comparison of unconditional pairs,
range-aware centered bounds, certified empirical minima, and estimated-mass
stratified pairs. That comparison uses the same eight scalar finite-support
laws, 300 repetitions at each $n\in\{64,256,1024,4096\}$, and separate
$0.05$ one-sided allocations per method; it does not assert one joint event
across methods. The supplement retains every population optimum, endpoint,
tolerance decision, and timing result.

On partial collapse at $n=4096$, stratification raises the population witness
from $0.0867$ to the exact risk $0.2533$. Mean calibrated endpoints are
$0.0646$ for unconditional pairs, $0.1772$ for estimated-mass stratification,
and $0.2178$ for empirical minimization. Finer conditioning is not uniformly
better: the chosen partition lowers the three-point population witness from
$0.2222$ to $0.1667$. Rare conflicts remain difficult. At $n=4096$, the
exact MSE is $0.025$, yet neither primary pair method excludes RMSE $0.10$
in any of 300 trials, whereas empirical minimization does so in every trial.
Supplying a tighter known score range raises stratified exclusions to 284;
also supplying known masses raises them to 300. The result depends on both
the pairing construction and the structural information provided.

All four methods have zero observed coverage violations in this prescribed
grid, which is a numerical check rather than a proof. Finite-support lookup
and repeated-count optimization make its computation times unsuitable as
continuous-representation scalability evidence. A separate two-atom
calibration check has one pair-lower violation in 28,800 trials and none for
the centered lower or fitted upper bound. Accounting for known zero
representation variance substantially strengthens centered calibration in
the collapsed case. Thus these experiments do not establish uniform
superiority of pairwise certificates over centered or empirical-minimum
alternatives.

\subsection{An analytic check of the cost decision}
\label{sec:decision-value}

An explicit-feature control uses independent histories $X=(Z,S)$ and targets
$Y=0.8SZ+0.1N$ (repairable) or $Y=0.8SR+0.1N$ (irreducible), with
$S\sim\operatorname{Bernoulli}(0.2)$ and independent random signs $Z,R,N$.
Its three retained classes have costs $1,3,5$ and exact minimum MSEs
$0.138,0.138,0.01$ in the repairable regime, and $0.138$ for all three
in the control. These are explicit maps, not learned KAHMs.
At RMSE tolerance $0.30$, upper calibration accepts the cost-5 fit in every
one of 100 repairable repetitions from $n=1024$. Stratified lower calibration
certifies its minimum feasible cost in every repetition from $n=16384$;
unconditional calibration never closes the cost interval through $n=65536$.
Both selection rules return the same fitted model: the additional conclusion
is necessity of its declared cost. The irreducible control is jointly excluded
in every repetition only at $n=65536$ under stratification, with smaller
samples and unconditional bounds retained as failures. The supplement gives
the complete design, all intermediate outcomes, and simultaneous allocation.
